\documentclass[11pt]{article}
\usepackage[margin=1in]{geometry}
\usepackage{amsmath,amssymb,amsthm}
\newtheorem{theorem}{Theorem}
\newtheorem{lemma}[theorem]{Lemma}
\newtheorem{corollary}[theorem]{Corollary}
\theoremstyle{definition}
\newtheorem{remark}[theorem]{Remark}
\DeclareMathOperator{\Imag}{Im}
\DeclareMathOperator{\Real}{Re}
\newcommand{\bb}{\mathbf}
\title{The two-row $d=4$ fiber law: reduction to a no-rational-root statement,\\
with a $2$-adic tower}
\author{Clio}
\date{6 June 2026}

\begin{document}
\maketitle

\begin{abstract}
For $\lambda=(2m-b,b)$ with $\psi=h_2+ie_2$, the $d=4$ fiber-vanishing law asserts
$G_\lambda:=\langle s_\lambda,\psi^m\rangle=0\iff\lambda=(2,2)$. Prior work reduced the two-row
case to $\Imag G_{(2m-b,b)}\neq0$ and closed $b\le4$. We give a clean reduction
$\Imag\big((1+su+u^2)^m\big)=u\,H_m(u)$ with $H_m\in\mathbb Z[u]$, prove that
$I_b(m):=\Imag G_{(2m-b,b)}$ is a polynomial of degree $b$ (or $b-1$ if $4\mid b$) with forced
integer roots exactly $\{0,\dots,\lfloor(b-1)/2\rfloor\}$, and thereby reduce the law to:
the explicit ``tail'' polynomial $Q_b$ has no integer root $\ge b$. Computation through
$b=24$ shows $Q_b$ is irreducible over $\mathbb Q$ (or a half-integer linear factor times an
irreducible when $4\mid b$), so the law is equivalent to the no-rational-root statement
$(\diamond)$. We prove the first two levels of a $2$-adic tower yielding explicit infinite
families of non-vanishing, and show $v_2(I_b(m))$ is unbounded with flat $2$-adic Newton
polygon---explaining why finite $2$-adic and Eisenstein-type certificates cannot close
$(\diamond)$.
\end{abstract}

\section{Setup and the engine identity}
Write $s=1+i$, $P(u)=(1+su+u^2)^m$, and $G_{(2m-b,b)}=[u^b]\big((1-u)P\big)$. Set
$I_b(m):=\Imag G_{(2m-b,b)}\in\mathbb Z$. Since $s-1=i$, with $W:=1+u+u^2$,
\[
  1+su+u^2=(1+u+u^2)+iu=W+iu=:A,\qquad B:=\bar A=W-iu,
\]
so that $A+B=2W$ and $AB=W^2+u^2$.

\section{The clean reduction}
\begin{theorem}\label{thm:clean}
$\Imag(A^m)=u\,H_m(u)$, where
$H_m(u)=\dfrac{A^m-B^m}{A-B}=\sum_{j=0}^{m-1}A^jB^{m-1-j}\in\mathbb Z[u]$
is symmetric in $A,B$, hence a polynomial in $2W$ and $W^2+u^2$. Consequently
\[
  I_b(m)=[u^{b-1}]\big((1-u)\,H_m(u)\big).
\]
\end{theorem}
\begin{proof}
As $B=\bar A$, $A^m-B^m=2i\,\Imag(A^m)$ and $A-B=2iu$, so $H_m=\Imag(A^m)/u$. Integrality is
the fundamental theorem of symmetric functions applied to the roots $A,B$ of
$t^2-(2W)t+(W^2+u^2)$. Finally $I_b(m)=[u^b]\big((1-u)\,\Imag P\big)
=[u^b]\big((1-u)\,u\,H_m\big)=[u^{b-1}]\big((1-u)H_m\big)$.
\end{proof}
The recurrence $h_n=2W\,h_{n-1}-(W^2+u^2)h_{n-2}$ ($h_0=1,h_1=2W$, $H_m=h_{m-1}$) gives
$H_m=(W^2+u^2)^{(m-1)/2}\,U_{m-1}\!\big(W/\sqrt{W^2+u^2}\big)$ with $U$ the Chebyshev
polynomial of the second kind.

\section{Degree and the forced integer roots}
\begin{theorem}\label{thm:roots}
\textup{(a)} $\deg_u(\Imag P)=2m-1$ with leading coefficient $m$; hence
$\deg_u\big((1-u)\Imag P\big)=2m$.\\
\textup{(b)} $I_b(m)=0$ for every integer $m$ with $2m<b$, i.e.\ $m\in\{0,1,\dots,\lfloor(b-1)/2\rfloor\}$.\\
\textup{(c)} As a polynomial in $m$, $\deg_m I_b=b$ if $b\not\equiv0\ (\mathrm{mod}\ 4)$ and
$b-1$ if $4\mid b$; in particular $I_b$ is a nonzero polynomial.
\end{theorem}
\begin{proof}
(a) $\Imag P=\sum_{k\ \mathrm{odd}}(-1)^{(k-1)/2}\binom mk u^kW^{m-k}$; the $k$-th term has
degree $2m-k$, maximal at $k=1$ with leading coefficient $\binom m1=m$.
(b) $(1-u)\Imag P$ has degree $2m$, so $[u^b]=0$ once $b>2m$.
(c) For fixed $l$, $[u^l]P=\sum_k\binom mk i^k T(m-k,l-k)$ with $T(N,j)=[u^j]W^N$. Since
$\binom mk\sim m^k/k!$ and $T(m-k,l-k)\sim m^{l-k}/(l-k)!$, the $k$-th summand is
$\sim\frac{i^k}{k!(l-k)!}m^l$, whence $[u^l]P\sim\frac{(1+i)^l}{l!}m^l$ and
$\Imag[u^l]P\sim\frac{2^{l/2}\sin(\pi l/4)}{l!}m^l$. This $m^l$-coefficient vanishes iff
$4\mid l$; applying it to $I_b=\Imag[u^b]P-\Imag[u^{b-1}]P$ gives the stated degree.
\end{proof}
Thus $I_b(m)=\Big(\prod_{r=0}^{\lfloor(b-1)/2\rfloor}(m-r)\Big)\,Q_b(m)$ with
$\deg Q_b=\lfloor b/2\rfloor$ (resp.\ $\lfloor b/2\rfloor-1$ if $4\mid b$). For $m\ge b$ the
product is nonzero, so the two-row law is equivalent to
\[
  (\star)\qquad Q_b(m)\neq0\ \text{for all integers } m\ge b.
\]

\section{Reduction to ``no rational root''}
\begin{theorem}[verified $b\le24$]\label{thm:irred}
If $b\not\equiv0\ (\mathrm{mod}\ 4)$, $Q_b$ is irreducible over $\mathbb Q$. If $4\mid b$,
$Q_b=(2m-(2b-1))\,R_b(m)$ with $R_b$ irreducible of degree $\lfloor b/2\rfloor-1$; the linear
factor gives the simple rational root $m=(2b-1)/2$, a half-integer. In all cases $Q_b$ has no
integer root, so the two-row law holds for $b\le24$.
\end{theorem}
The factor $2m-(2b-1)$ for $4\mid b$ reflects $I_b\big((2b-1)/2\big)=0$, verified directly for
$b=4,8,\dots,28$. Hence the entire remaining gap is the single assertion
\[
  (\diamond)\qquad\text{for }b\ge5,\ Q_b\text{ has no rational root, except the simple }
  (2b-1)/2\text{ when }4\mid b.
\]

\paragraph{Why cheap certificates fail.} For the primitive integer model of $Q_b$, both the
constant and leading coefficients are odd: the $2$-adic Newton polygon is the flat slope-$0$
segment, so all roots are $2$-adic units (no $2$-adic ramification, no Eisenstein prime at
$2$). And since $\deg Q_b\to\infty$, no fixed prime $p$ can give ``no root mod $p$''
uniformly. So $(\diamond)$ requires genuine arithmetic input.

\section{The $2$-adic tower}
Mod $2$, $W^2+u^2\equiv 1+u^4$.
\begin{theorem}\label{thm:tower}
Using $H_m=\sum_c(-1)^c\binom{m-1-c}{c}(2W)^{m-1-2c}(W^2+u^2)^c$:

\emph{Level 1.} If $m$ is even, all coefficients of $H_m$ are even, so $I_b(m)$ is even. If
$m$ is odd, $H_m\equiv(1+u^4)^{(m-1)/2}\pmod2$, whence with $M=(m-1)/2$,
\[
  I_b(m)\equiv
  \begin{cases}
    \binom{M}{(b-1)/4} & b\equiv1\\
    \binom{M}{(b-2)/4} & b\equiv2\\
    0 & b\equiv0,3
  \end{cases}\pmod2\quad(\mathrm{mod}\ 4\ \text{on }b).
\]

\emph{Level 2.} If $m\equiv0\pmod2$, $M'=(m-2)/2$, then $H_m\equiv 2(-1)^{M'}(m/2)\,W\,(1+u^4)^{M'}\pmod4$,
and since $(1-u)W=1-u^3$,
\[
  \tfrac12 I_b(m)\equiv
  \begin{cases}
    (m/2)\binom{M'}{(b-1)/4} & b\equiv1\\
    (m/2)\binom{M'}{(b-4)/4} & b\equiv0\\
    0 & b\equiv2,3
  \end{cases}\pmod2.
\]
\end{theorem}
\begin{corollary}
$I_b(m)\neq0$ whenever ($m$ odd, $b\equiv1,2$, $\binom{(m-1)/2}{\lfloor(b-1)/4\rfloor}$ odd)
or ($m\equiv2\ (\mathrm{mod}\ 4)$, $b\equiv0,1$, level-2 binomial odd). By Lucas these are
infinite families for each admissible $b$.
\end{corollary}
\begin{remark}
$v_2(I_b(m))$ is unbounded (it reaches $11$ at $(m,b)=(18,6),(18,10)$), so no finite
truncation of the tower closes $(\diamond)$: $m$ can be $2$-adically arbitrarily close to a
($2$-adic unit) root of $Q_b$, forcing arbitrarily large $v_2$---yet never equal, as those
roots are irrational. This is the $2$-adic shadow of $(\diamond)$.
\end{remark}

\section{Status}
Theorems \ref{thm:clean}, \ref{thm:roots}, \ref{thm:tower} are rigorous;
Theorem \ref{thm:irred} is verified through $b=24$. The full law is equivalent to
$(\diamond)$, an irreducibility/no-rational-root statement about an explicit integer-valued
polynomial family of degree $\lfloor b/2\rfloor$ with flat $2$-adic Newton polygon---likely a
disguised classical orthogonal family, whose identification is the natural next step.
\end{document}
